\documentclass[11pt]{article}
\usepackage{manual_docs_style}
\externaldocument{build/environment_profiles}[environment_profiles.pdf]
\externaldocument{build/generation_policy_schema}[generation_policy_schema.pdf]
\externaldocument{build/run_graph_schema}[run_graph_schema.pdf]

\begin{document}

\docTitle{Stage: Candidate Run Generation}
\phantomsection\label{docstart:candidate_generation_stage}

\section*{Overview}

This stage turns a high-level sampling policy into a resolved run graph.
It is the only stage that samples new baseline parameters, intervention times, intervention values, and candidate families.
Stage: Simulate does not sample; it only materializes the resolved recipes produced here.

The main driver is:

\begin{DocCode}
workflows/generate_runs/compile_run_graph.py
\end{DocCode}

A typical invocation is:

\begin{DocCode}
python workflows/generate_runs/compile_run_graph.py \
  --model BallDrop \
  --policy benchmark_v1 \
  --overwrite
\end{DocCode}

\section*{Inputs}

For each model under \code{models/simulink/<MODEL>/}, the stage reads:

\begin{itemize}
\item \code{experiment_config.json}: hard environment contract, including exposed variables, legal intervals, simulation horizon, observed signals, and detectability configuration.
\item \code{generated/metadata.json}: generated Simulink metadata, including exposed workspace variables and intervention block mappings.
\item \code{generation_policies/<policy_id>.json}: campaign-specific sampling recipe.
\item \code{manual_cases.jsonl}: optional curated cases to merge into the resolved run graph.
\end{itemize}

\section*{Outputs}

The stage writes a plan under:

\begin{DocCode}
models/simulink/<MODEL>/plans/<policy_id>/
\end{DocCode}

with these artifacts:

\begin{itemize}
\item \code{run_nodes.jsonl}: flat list of runnable baseline, intervention, and time-zero/static-change recipes.
\item \code{run_edges.jsonl}: graph relationships among run nodes.
\item \code{generation_report.json}: generation summary, coverage statistics, and skipped parameter-direction requests.
\item \code{validation_report.json}: schema, semantic, and policy-acceptance validation summary.
\end{itemize}

\section*{Execution Flow}

\begin{enumerate}
\item Load \code{experiment_config.json} and validate that all exposed variables are present in \code{generated/metadata.json}.
\item Load \code{generation_policies/<policy_id>.json}.
\item Resolve the random seed and initialize deterministic random-number generators.
\item Sample baseline families according to \code{baseline_sampler}.
\item For each baseline family, sample intervention children according to \code{intervention_sampler}.
\item For each intervention child, create a matched time-zero/static-change recipe.
\item Add optional curated manual cases from \code{manual_cases.jsonl}.
\item Canonicalize every recipe and compute \code{recipe_hash} and \code{run_id}.
\item Write flat runnable nodes to \code{run_nodes.jsonl}.
\item Write baseline/intervention/time-zero relationships to \code{run_edges.jsonl}.
\item Write skipped or impossible parameter-direction requests to \code{generation_report.json}.
\item Validate the resolved run graph and write \code{validation_report.json}.
\end{enumerate}

\section*{Policy Intent and Acceptance Criteria}

A generation policy should state the intended use of the candidate pool it creates.
This is not a scientific hypothesis about agent behavior; it is an engineering contract for deciding whether the generated candidate universe is sufficient for the next stages.

A generated plan is accepted only if it satisfies the policy's acceptance criteria.
Typical criteria include:

\begin{itemize}
\item all generated run nodes pass schema validation;
\item all \code{run_id} and \code{recipe_hash} values are unique;
\item all sampled values lie inside \code{experiment_config.json} legal intervals;
\item every intervention changes exactly one root-cause parameter;
\item every intervention has a matched no-intervention baseline and time-zero/static-change run;
\item the number of generated candidate families is at least the requested target;
\item skipped parameter-direction requests are recorded in \code{generation_report.json};
\item policy coverage targets are either satisfied or explicitly reported as unmet.
\end{itemize}

The candidate-generation stage does not guarantee that intervention trajectories are detectable or distinguishable after simulation.
Those properties are evaluated by the metric and oracle stages after materialization.

\section*{Manual Cases}

Manual cases are optional curated examples.
They are not a separate downstream schema.
The generator resolves them into the same \code{run_nodes.jsonl} and \code{run_edges.jsonl} format as programmatic cases.

Manual cases must still satisfy core validity rules: finite values, legal intervals, valid intervention times, and exactly one changed root-cause parameter for intervention runs.
Manual provenance is recorded in node metadata:

\begin{DocCode}
"metadata": {
  "source": "manual",
  "validation_profile": "diagnostic"
}
\end{DocCode}

\section*{Skipped Directions}

The resolved run graph contains only runnable recipes.
If the generator cannot sample an increase or decrease side for a parameter because the legal interval and distance constraints make it impossible, it must not create a null-valued run node.
Instead, it records the skipped request in \code{generation_report.json}.

\section*{Artifact Policy}

\begin{itemize}
\item \code{generation_policies/<policy_id>.json} is a source artifact and may be human-authored.
\item \code{run_nodes.jsonl}, \code{run_edges.jsonl}, \code{generation_report.json}, and \code{validation_report.json} are generated artifacts.
\item To change a generated candidate pool, edit the policy or manual cases and regenerate the plan.
\item New downstream code should consume the resolved run graph, not \code{model_run_specs.json}.
\item A legacy \code{model_run_specs.json} compatibility export may be generated during migration, but it is not canonical.
\end{itemize}

\end{document}
